\numberlesssection{РЕФЕРАТ}
Расчётно-пояснительная записка \pageref{LastPage} страниц, \total{figure} рисунков,
\total{table} таблицы, 10 источников.

\noindent
КОМПЬЮТЕРНАЯ ГРАФИКА, КИНЕМАТИКА, ШЕСТЕРЕНОЧНЫЙ МЕХАНИЗМ, АЛГОРИТМ Z-БУФЕРА, ЗАКРАСКА ФОНГА, ПЛАНАРНЫЕ ТЕНИ, ГЕКСАГОНАЛЬНАЯ СЕТКА, C++, QT.

Объект исследования --- методы визуализации трехмерных сцен и алгоритмы моделирования кинематических связей в механизмах.

Цель работы --- разработка программного обеспечения для интерактивного конструирования и визуализации кинематики шестереночных механизмов в трехмерном пространстве.

В работе использованы методы аналитической геометрии и линейной алгебры для описания трехмерных объектов и их преобразований. Для визуализации применен метод программной растеризации с использованием алгоритма Z-буфера для удаления невидимых поверхностей. Расчет освещения выполнен с использованием модели Фонга. Для построения теней использован метод планарной проекции. Расчет кинематики базируется на представлении механизма в виде графа и алгоритме поиска в ширину (BFS).

Разработано программное обеспечение, позволяющее конструировать механизмы из зубчатых колес на гексагональной сетке, визуализировать их вращение с учетом физических связей и определять заклинивание. Реализовано корректное отображение освещения и теней.

Проведен анализ производительности системы. Установлено, что время отрисовки кадра линейно зависит от количества объектов в сцене. Программа обеспечивает интерактивную частоту кадров (около 18 FPS) при отрисовке 100 шестеренок.
